% !TeX root = ../papers/nc-p3.tex
\section{偶交換を外す楕円曲線と最大反対階数}
\label{sec:klein-elliptic-completion}
四位置の文字空間を独立な四元群の正則表現とし、テンソルの総文字次数を零とする。
偶交換を仮定しない。前節までの片側対称点の近くにある二つの隣接切断を
\[
 M_h(i,k)=w_{i,i+h,k,k+h},\qquad
 N_h(i,k)=w_{k,i+h,i,k+h}=M_{h+i+k}(k,i)
\]
と書く。$M_h,N_h$はいずれも$h=0,1$で階数2、$h=2,3$で階数1とする。

\begin{theorem}\label{thm:klein-elliptic-completion}
二つの隣接切断に上の文字次元を持つ窓には、反対階数16、最小周期4で
幾何学的一般点が標準Koszul左右AS正則となる既約な族がある。
印を忘れた添字保持同型類の像は五次元である。その標準形の最後の一変数は
一般の四つの係数に対して滑らかな種数1曲線を動く。
この族は、以下の非退化な表示の一般の基底係数の上にある全解を覆う。
\end{theorem}

\paragraph{二切断の同時規格化。}
二つの薄い色の全成分が非零とする。$M_2,N_2$は対角が等しい階数1行列なので、
偶二位置の相対的な対角基底変更で$M_2=N_2$とできる。この後の偶倍率は
両位置で射影的に同じである。$N_3$も$M_3$の対角共役であり、
各組$A=\{0,1\},B=\{2,3\}$内の非対角成分は
$N_2=M_2$から等しい。従って共役の対角倍率は$A$上で$\lambda$、$B$上で1と取れる。

さらに$M_h(i,k)=u_i^h v_k^h$と書けば、$M_2(i,\bar i)=M_3(\bar i,i)$から
$u_i^2v_{\bar i}^2=u_{\bar i}^3v_i^3$となる（上付きは色を表す）。
よって奇位置の各列の倍率で、両色について$v_k^h=u_k^{h+1}$にできる。
同じ奇列に現れる二つの条件の整合性がまさにこの等式である。
薄い係数は以前と同じ二行列$U,V$で表され、前と同じ倍率変更で
$U$を全成分1、$V$を右下$z$以外1とできる。

$A_0=\left(\begin{smallmatrix}a&b\\c&d\end{smallmatrix}\right)$,
$A_1=\left(\begin{smallmatrix}e&f\\g&h\end{smallmatrix}\right)$を$M_0,M_1^{\mathsf t}$の
左上とする。残る倍率で$A_0$の共通スカラーを除ける。
$X_l=\left(\begin{smallmatrix}1&1/l\\1/l&z\end{smallmatrix}\right)$,
$Y_l=\left(\begin{smallmatrix}1&l\\lz&1\end{smallmatrix}\right)$,
$F_l(A)=Y_l A^{-1}X_l$と置く。
$M_0$の交差ブロックは$X_\lambda,Y_\lambda$、
$M_1^{\mathsf t}$では$\lambda X_\lambda,Y_\lambda/\lambda$となる。
これは行と列が異なる組にあるとき、$N$の薄い側で
色3の成分だけが組の向きによって$\lambda^{\pm1}$倍されるためである。
従って両者の右下は$T=F_\lambda(A_0),T'=F_\lambda(A_1)$である。
$N_0,N_1^{\mathsf t}$の左上は$C=\left(\begin{smallmatrix}a&f\\g&d\end{smallmatrix}\right)$,
$D=\left(\begin{smallmatrix}e&b\\c&h\end{smallmatrix}\right)$となる。
これらの右下は$F_1(C),F_1(D)$でなければならず、同時にそれぞれ
$T$の対角と$T'$の非対角、$T'$の対角と$T$の非対角を組み合わせた行列である。
以後、使う各二行列の行列式と表示の分母が非零な部分で考える。

\paragraph{一つの判別式を持つ二次方程式。}
$x=\det C=ad-fg$, $y=\det D=eh-bc$とする。
$F_1(C)$の対角二式と$F_1(D)$の非対角二式は独立に解け、
\[
 f=f_0+f_1x,\quad g=g_0+g_1x,\qquad e=e_0+e_1y,\quad h=h_0+h_1y
\]
となる。具体的には
\[
\begin{split}
 f_0&=(a+d)/(z+1),\quad g_0=z f_0,\quad
 f_1=(T_{00}-T_{11})/(z^2-1),\quad g_1=-T_{00}-f_1,\\
 e_0&=h_0=(zb+c)/(z+1),\quad
 e_1=(zT_{01}-T_{10})/(z^2-1),\quad h_1=(T_{10}-e_1)/z .
\end{split}
\]
$P(x)=fg+x-ad=p x^2+m x+n$、$Q(y)=eh-y-bc=q y^2+k y+o$とする。
両者の判別式は等しく、$\delta=m^2-4pn=k^2-4qo$である。
二根を
\[
             2p x+m=s,\quad 2q y+k=-s,\qquad s^2=\delta
\]
と組ませると、残る四つのSchur補の等式も全て従う。
これらは六つの独立な係数$a,b,c,d,z,\lambda$の有理関数体上の恒等式である。
\path{computations/verify_klein_elliptic_completion.py}は、第一に二つの判別式を引き、
第二に$y=-(m+k)/(2q)-(p/q)x$を入れ、
残る等式の分母を払って$P(x)$で割った余りの定数項・一次項を照合する。
$\det A_1=x+y-ad+bc$なので二次拡大での行列の逆計算は不要である。

$\lambda^2\delta$は$\lambda$の高々四次の多項式である。実際、$f_1,g_1$は
$\lambda,1,\lambda^{-1}$の線型結合である。$(a,b,c,d,z)=(5,2,3,7,-2)$では
この多項式の7569倍は
\[
 31024\lambda^4-17712\lambda^3-22355\lambda^2+8856\lambda+7756
\]
で、四根は相異なる。一般の基底係数でも四つの単純分岐点があり、滑らかな射影モデル
の種数は1である。二重被覆は幾何学的に既約である。
もう一方の符号による組合せは一般には残る四式を満たさない。
例えば同じ五係数で$\lambda=2$とし、$2q y+k=2p x+m$を代入すると
対角の最初の等式の余りは非零になる（同コード）。$P$が一般に既約なので、
この組合せの解が一般の基底へ支配的に写ることはない。

\paragraph{正則性と同型類の次元。}
$\lambda=1$では$f=b,g=c,x=ad-bc$の有理切断があり、
上記の符号で$(e,h)$は片側対称族の第二補完の対角を与える。
実際、この切断における$P,Q$の符号関係は二つの有理関数の平方の一致である。
$(5,2,3,7,-2)$で$s=1$, $(e,h)=(1/6,-1/4)$となることで符号が決まる。
従って四次中心により正則と証明した点がこの既約族上にある。有限Frobenius性の
開条件と内部次数別のbar複体を用いた従来の特殊化論法により、幾何学的一般点も標準AS正則である。
反対階数16は下の例で達するため一般値である。
階数16では、どの二位置の正または負の交換が反対対にある場合も、その対の
像次元を10または6に抑えるので不可能である。隣接対の交換が可能なら残る二分割の
切断階数が等しくなるが、6と16なのでこれも不可能である。

片側対称切断の一般点で既知の被約二十点配置は、どの三点も共線でなく、
五点平面も持たない。後者は節\ref{sec:klein-generic-isomorphisms}の特殊点での
全平面の計算と有限エタール順序付け被覆から従う。
同節の射影行列の固有な比較を各辺に適用でき、両奇配置の非同値性が近傍で
保たれる。従って一般の最小周期は4である。また各点配置の射影安定群は有限なので、
四元群の印を忘れるファイバーは有限となる。上の規格化後は、
薄い次数の二つの積比から$z$が決まり、$N_3/M_3$の対角共役比から$\lambda$が決まる。
$M_2=N_2$が偶二位置の相対倍率を、$M_2=M_3^{\mathsf t}$が対応する奇二位置の
相対倍率を固定する。文字を固定した残る対角作用は前節までと同じで、
$A_0$の共通スカラーを除けば有限である。従って基底の四変数と楕円曲線の
一変数を合わせ、印なしの同型類の像は五次元となる。

\begin{proposition}\label{prop:klein-rational-full-opposite}
$(a,b,c,d,z,\lambda)=(1,1,1,2,-2,-2)$、$s=-27/2$から得る有理窓は、
反対階数16の標準Koszul左右AS正則窓であり、四つの中心四次式上で各始点が自由階数256となる。
\end{proposition}
\begin{proof}
$(e,f,g,h)=(-13/14,1/13,-14,-16/11)$となり、反対切断の行列式は
$-29506135761/19712$である。隣接切断はともに階数6、一文字切断は反対切断が
可逆なので全階数である。四辺の関係空間は$M,N$の列空間と行空間から読み、
奇始点では二文字を反転する。従って片側対称な場合の反復関係を仮定する必要はない。
実際の四窓について中心計算と商計算を再実行すると、文字次数順の中心次元は
$(3,2,0,0)$で、定理\ref{thm:klein-rational-period-four}と同じ選び方で
四式を割った商は全始点で次数13まで$(1+t+t^2+t^3)^4$の係数以下である。
幾何学的一般点の標準Hilbert関数が下界となるため、先と同じ自由加群と
Koszul列の双対の議論で全次数の主張が従う。
\end{proof}

\begin{corollary}\label{cor:klein-four-configurations}
上の反対階数16の有理点自身で、四位置の被約二十点配置は全て互いに射影非同値で、
それぞれの射影安定群は四元群である。従って最小周期は4であり、
両方の反対対で二次Veronese表示が相異なる。これらは楕円曲線族の一般点にも成立する。
\end{corollary}
\begin{proof}
素数2657で四辺それぞれの最大小行列式を四つの射影座標層で消去する。
全四位置で幾何点は相異なる素体上の二十点であり、核を四回取ると元の番号に戻る。
期待余次元の長さ20から各配置は被約である。全四配置で三点直線はなく、
四点平面48枚のうち八枚を通す四点が線型に独立である。
それらと第五点の像から写像を全列挙した個数行列は$4 I_4$となる。
\path{computations/verify_klein_full_rank_configurations.py}で検算する。
節\ref{sec:klein-generic-isomorphisms}と同じ固有な行列比較により、これが標数零での上界となる。
指定された群が全配置を保つため等号である。
定理\ref{thm:finite-square-chart}により反対位置の平方空間は単元倍で対応するが、
同じVeronese三次元多様体に載る二つの表示なら、その二十点配置は射影同値になってしまう。
\end{proof}
